% =====================================================================
%  大学別オンライン数学模試 ── 返却見本「合格への手引き」
%
%  1枚目は成績表。得点・分野別・大問別をグラフで出し、1枚で全体が分かるようにする。
%  2枚目以降は答案への講評で、どこで何を取り違えたかを答案に即して書く。
%
%  組版は「合格答案をつくる」シリーズの本文を土台にするが、
%  成績表だけは色を使う。白黒前提の本文と違い、これは画面と手元で読むもので、
%  得点の多い少ないを一目で掴ませたいため。
%
%  中の受験者・得点・答案はすべて架空のもの。実在の受験者のものではない。
%  問題は見本問題（数列）と同じもので、答えは
%  scripts/verify-solutions.py で確かめてある。
% =====================================================================
\input{preamble}

\renewcommand{\seriesname}{大学別オンライン数学模試}
\renewcommand{\booktitle}{合格への手引き}
\renewcommand{\bookyear}{2027}

% ---- 成績表の配色。サイトの色をそのまま持ってくる ----
\definecolor{mTeal}{HTML}{1F5F5B}   % 模試の色
\definecolor{mNavy}{HTML}{1B3A63}   % 得点の棒
\definecolor{mRed}{HTML}{B3412F}    % 落としたところ
\definecolor{mAmber}{HTML}{8A5A1C}  % あと一歩
\definecolor{mGreen}{HTML}{2F6043}  % 取れているところ
\definecolor{mRule}{HTML}{E0E3E7}
\definecolor{mPale}{HTML}{F2F6F5}

\usepackage{colortbl}

% 見出しの帯
\newcommand{\cardhd}[1]{%
  \par\noindent\colorbox{mTeal}{\parbox{\dimexpr\linewidth-2\fboxsep\relax}{%
    \rule{0pt}{2.6ex}\hspace{0.2em}\textcolor{white}{\bfseries\small #1}\rule[-0.9ex]{0pt}{0.9ex}}}%
  \par\vspace{1.6mm}}

% 横棒。得点 / 満点 を色で描く
\newcommand{\scorebar}[4]{% #1 ラベル #2 得点 #3 満点 #4 色
  \noindent\makebox[7.2em][l]{\small #1}%
  \begin{tikzpicture}[baseline=-0.6ex]
    \fill[mRule] (0,0) rectangle (6.4,0.30);
    \fill[#4] (0,0) rectangle ({6.4*#2/#3},0.30);
  \end{tikzpicture}%
  \hspace{0.6em}{\small\bfseries #2}{\footnotesize\textcolor{black!55}{\,/\,#3}}\par\vspace{1.3mm}}

\begin{document}

% =====================================================================
%  1枚目：成績表
% =====================================================================
\newgeometry{top=14mm,bottom=12mm,left=16mm,right=16mm}
\thispagestyle{empty}
\begingroup
\renewcommand{\baselinestretch}{1.12}\selectfont
\setlength{\parskip}{0pt}
\begin{center}
  {\footnotesize\textcolor{mTeal}{\bfseries 2027年度入試向け ・ 大学別オンライン数学模試}}\\[1.2mm]
  {\fontsize{17}{20}\selectfont\bfseries 成\hspace{0.4em}績\hspace{0.4em}表}\\[1.0mm]
  {\footnotesize （返却見本・記載内容はすべて架空のものです）}
\end{center}
\vspace{2.4mm}

\noindent\colorbox{mPale}{\parbox{\dimexpr\linewidth-2\fboxsep\relax}{%
\vspace{1.2mm}
\hspace{0.4em}\begin{minipage}{0.97\linewidth}\small
\begin{tabular}{@{}p{4.6em}p{12.5em}p{6.4em}p{12em}@{}}
受験者 & \textbf{（見本）受験 太郎} & 受験した模試 & \textbf{三重大学 理系数学} \\
受験日 & 2026年12月3日 & 返却日 & 2026年12月15日 \\
\end{tabular}
\end{minipage}
\vspace{1.2mm}}}

\vspace{3mm}
\cardhd{1　得点}

\noindent
\begin{minipage}[t]{0.46\linewidth}
  \vspace{0pt}
  \begin{center}
    {\fontsize{34}{38}\selectfont\bfseries\textcolor{mNavy}{13}}%
    {\large\textcolor{black!55}{\,/\,20}}\\[1.0mm]
    {\footnotesize 得点率 \textbf{65\%}}
  \end{center}
\end{minipage}%
\begin{minipage}[t]{0.52\linewidth}
  \vspace{2mm}
  \scorebar{大問1（数列）}{13}{20}{mNavy}
  \vspace{2mm}
  \noindent{\footnotesize\textcolor{black!55}{%
  この見本は大問1題のみです。実際の模試は大問3〜6題で実施します。}}
\end{minipage}

\vspace{3.4mm}
\cardhd{2　小問ごとの得点}

\noindent
\scorebar{(1) 初項と第3項}{4}{4}{mGreen}
\scorebar{(2) 漸化式・一般項}{6}{8}{mAmber}
\scorebar{(3) $\sum k\,a_k$ の計算}{3}{8}{mRed}

\vspace{1.5mm}
\noindent{\footnotesize
\textcolor{mGreen}{■}\,満点\quad
\textcolor{mAmber}{■}\,あと一歩\quad
\textcolor{mRed}{■}\,立て直しが要る}

\vspace{3.4mm}
\cardhd{3　何で落としたか}

\noindent\small
\begin{tabular}{@{}p{3.2em}p{12em}p{18.5em}@{}}
\toprule
\textbf{小問} & \textbf{落とした理由} & \textbf{次にすること} \\
\midrule
(2) & 等比数列の初項の取り違え & 置き換えた数列に名前を付け、$b_1$ を置き換えた式から出す \\
\addlinespace[2pt]
(3) & ずらして引くときの両端の扱い & 2つの和を縦に並べて書いてから引く \\
\bottomrule
\end{tabular}

\vspace{2mm}
\noindent{\footnotesize\textcolor{black!60}{%
どちらも方針は正しく立てられています。落としているのは手の動きで、
同じ操作を書き方を決めて繰り返せば減ります。くわしくは次ページの講評を。}}

\vspace{3.4mm}
\cardhd{4　受験者全体の中での位置}

\noindent\colorbox{black!4}{\parbox{\dimexpr\linewidth-2\fboxsep\relax}{%
\vspace{1.2mm}\hspace{0.4em}\begin{minipage}{0.96\linewidth}\small
\textbf{この回は非掲載です。}\par\vspace{0.8mm}
平均点・偏差値・順位分布は、その大学の受験者が\textbf{30名以上}になった回から掲載します。
少人数の中での位置を数字にしても、意味のある値にならないためです。
掲載する回では、この欄に平均点・偏差値・得点分布のグラフが入ります。
\par\vspace{0.8mm}
合否判定は、どの回でも出していません。
\end{minipage}\vspace{1.2mm}}}

\vspace{3.4mm}
\cardhd{5　この2週間ですること}

\noindent\small
\begin{enumerate}[leftmargin=1.6em,itemsep=2pt,topsep=1pt]
  \item 置き換えの初項を、置き換えた式から出す練習を10題（$b_1$ だけでよい）
  \item $\dsp\sum_{k=1}^{n}k\,r^{\,k}$ を、縦に並べて引く書き方で5題
  \item 一般項を出したら、必ず最初の3項で確かめる（上の15題すべてで）
\end{enumerate}

\vfill
\noindent{\footnotesize\textcolor{black!55}{%
本模試は三重大学とは関係のない、当サイトが独自に制作・実施するものです。
配点・採点基準は本模試のもので、大学が公表しているものではありません。
この見本の受験者・得点・答案はすべて架空のものです。\hfill yuta-eng.com}}

\endgroup
\clearpage
\restoregeometry
\resetcolht

% =====================================================================
%  2枚目以降：答案への講評
% =====================================================================
\section*{答案への講評}

\hdA{(2) 漸化式と一般項 ── 6点 / 8点}

\hdB{書けていたところ}
$a_{n+1}=S_{n+1}-S_{n}$ を使うと最初に宣言し、差をとって
$a_{n+1}=2a_{n}+3$ まで正しく整理できています。ここまでで4点。
方針の宣言が1行あるだけで、読む側は以降の計算を追いやすくなります。この書き方は続けてください。

\hdB{落としたところ}
両辺に $3$ を足して $\{a_{n}+3\}$ が等比数列になることまでは書けていますが、
その初項を $a_{1}=3$ としています。等比数列になっているのは $a_{n}+3$ のほうなので、初項は
\[
  a_{1}+3=6
\]
です。ここを取り違えたため、一般項が $a_{n}=3\cdot2^{\,n-1}$ となり、
正しい $a_{n}=3\left(2^{\,n}-1\right)$ からずれました。

\hdB{次に直すこと}
\textbf{置き換えた数列の初項は、置き換えた式に $n=1$ を入れて出す。}
$\{b_{n}\}=\{a_{n}+3\}$ と名前を付けてしまえば、$b_{1}=a_{1}+3=6$ と機械的に書けます。
名前を付けずに進めると、元の数列の初項をそのまま使ってしまいがちです。

\smallskip
また、(1) で $a_{1},a_{2},a_{3}=3,9,21$ を出しているので、
一般項に $n=1,2,3$ を入れて突き合わせれば、この誤りはその場で見つかりました。
\textbf{一般項を出したら必ず最初の数項で確かめる}ことを習慣にしてください。

\hdA{(3) $\dsp\sum_{k=1}^{n}k\,a_{k}$ ── 3点 / 8点}

\hdB{書けていたところ}
$k\,a_{k}=3k\cdot2^{\,k}-3k$ と2つの和に分け、後ろの和を $\dfrac{n(n+1)}{2}$ と
処理するところまでは正しく書けています。方針としては最短のものを選べています。

\hdB{落としたところ}
$U_{n}=\dsp\sum_{k=1}^{n}k\cdot2^{\,k}$ をずらして引くところで、
$k=1$ の項と $k=n+1$ の項の扱いが抜けています。答案では
\[
  U_{n}-2U_{n}=\sum_{k=1}^{n}2^{\,k}
\]
としていますが、正しくは
\[
  U_{n}-2U_{n}=2+\sum_{k=2}^{n}2^{\,k}-n\cdot2^{\,n+1}
\]
で、先頭の $2$ と末尾の $-n\cdot2^{\,n+1}$ が落ちています。
ここが落ちると、以降の計算がすべてずれます。

\hdB{次に直すこと}
\textbf{ずらして引くときは、2つの和を縦に並べて書く。}
\[
  \begin{array}{rcccccc}
    U_{n}  &=& 1\cdot2 &+& 2\cdot2^{2} &+\cdots+& n\cdot2^{\,n}\\[1mm]
    2U_{n} &=&         & & 1\cdot2^{2} &+\cdots+& (n-1)\cdot2^{\,n}+n\cdot2^{\,n+1}
  \end{array}
\]
並べてから引けば、ずれて余る両端の項が目で見えます。
頭の中で引くと、必ずこの2か所を落とします。

\hdA{今後の学習の助言}

\hdB{三重大学の出題傾向と照らして}
当サイトの分析では、三重大学の理系数学は試験時間120分・大問3題・完全記述式で、
2019〜2026年度の8年分で三角関数・ベクトル・確率が各8題と多く、数列もこれに次ぎます。
大問3題のうち1題が数列になる年が続いているため、
\textbf{今回落とした2か所は、本番でそのまま出る可能性の高い操作}です。

\hdB{そのあとに}
この2点が固まったら、数列と他分野の融合（確率漸化式・数列と極限）に進んでください。
三重大学では分野をまたぐ出題が見られるため、単独の数列が解けるだけでは足りません。

\end{document}
